\documentclass[10pt]{article}

\usepackage[margin=0.85in]{geometry}
\usepackage{pgfplots}
\pgfplotsset{compat=1.18}
\usepackage{booktabs}
\usepackage{amsmath}
\usepackage{hyperref}
\hypersetup{hidelinks}
\usepackage{caption}
\usepackage{enumitem}
\usepackage{microtype}
\usepackage{placeins}

\captionsetup{font=small,labelfont=bf}
\setlength{\parskip}{0.35em}
\setlength{\parindent}{0pt}
\setlist[itemize]{leftmargin=1.4em,nosep}

\title{\vspace{-1.5em}COMP20003 Assignment 2: Stage 3 Report\\
\large Linked-List and Patricia-Trie Dictionary Complexity}
\author{Nhat Chi Nguyen (1492182)}
\date{}

\begin{document}
\maketitle
\vspace{-2.3em}

\section{Scope and theoretical expectations}

The Stage 1 dictionary (\texttt{dict1}) stores all $n$ records in a singly linked list.
Because duplicate keys are allowed and every matching record must be returned, this
implementation scans the entire list for every query. The Stage 2 dictionary
(\texttt{dict2}) stores keys in a binary Patricia trie. A search compares the query with
the packed bit-stem in each visited node and follows one branch until it reaches a leaf or
the first mismatching bit. Records with the same key are retained in input order at the
same leaf.

Both programs report bit comparisons ($b$), node accesses ($n_c$), and whole-string
comparisons ($s$). These counters measure different operations, so they should not be
treated as interchangeable proxies for running time. Let $N$ be the number of records,
$U$ the number of distinct keys, and $L$ the query length in bits, including the
\texttt{\textbackslash0} terminator. The expected lookup behaviour is:

\begin{table}[h]
\centering\small
\begin{tabular}{lll}
\toprule
Structure & Counter & Theoretical search cost \\
\midrule
Linked list & node/string & $\Theta(N)$ for this full-scan implementation \\
Linked list & bit & $O(NL)$ worst case; $\Theta(N)$ when $L$ is fixed \\
Patricia trie & bit & $\Theta(L)$ for a successful search; $O(L)$ for an unsuccessful one \\
Patricia trie & node & expected $O(\log U)$ for balanced branching; \\
& & worst case $O(\min\{U,L\})$ \\
\bottomrule
\end{tabular}
\caption{Predictions used to interpret the experiments.}
\label{tab:theory}
\end{table}

The closest-candidate edit-distance work is excluded from the reported counters, as
permitted by the specification. Consequently, the misspelling experiments measure only the
cost of locating the mismatching subtree; enumerating and ranking its candidates is a
separate, excluded cost. The report focuses on lookup rather than insertion.

\section{Experimental design}

A Python harness generated data, supplied queries to the submitted \texttt{dict1} and
\texttt{dict2} binaries, and parsed each program's \texttt{b\#\# n\#\# s\#\#} output.
Comparison counts are deterministic for a fixed dataset and query, so all variability
comes from which dataset and which queries were sampled, not from timing noise. To
capture both sources, each reported point in Experiments 1 and 2 pools results from 5
independently generated datasets, each queried with 20 samples (100 samples per point in
total), and reports the standard deviation across this pooled sample alongside the mean.
Both programs were first checked against Ed's automated test suite and the provided
reference outputs, so the experiments below measure comparison behaviour under
implementations that passed the available tests and matched the provided reference
outputs, rather than testing correctness itself.

\textbf{Experiment 1: growth and key density.}
Ten sizes from $N=50$ to $N=40\,000$ were tested with 32-character lowercase keys.
The \emph{sparse} datasets used independently generated keys, which usually diverged
within the first few bits. Each \emph{dense} key differed from an all-\texttt{a} baseline
at exactly two character positions, creating much longer pairwise agreement, with the
position pair and character pair drawn from a precomputed bijection so that no two keys
in a dataset could coincide. This experiment varies input size and the distribution of
informative bits.

\textbf{Experiment 2: edit position.}
At $N=5\,000$, three query types were compared at both densities: an exact stored key,
the same key with its last character changed (a \emph{last-character edit}), and the same
key with its first character changed (a \emph{first-character edit}). The replacement
character is drawn from the same alphabet as the stored keys, so an edit's position does
not necessarily coincide with where the trie's search first diverges from the query: the
modified prefix may still follow a branch shared with a different stored key, especially
under the dense generator's heavy prefix sharing.

All synthetic keys were unique, so $N=U$ in Experiments 1 and 2; this was asserted
programmatically for every generated dataset rather than assumed. Modified queries were
also checked to ensure they were not existing keys, redrawing the changed character when
a collision was found.

\textbf{Experiment 3: provided data.}
The provided \texttt{dataset\_1353.csv} was queried using \texttt{test1353.in} (588 exact
queries) and \texttt{testpart1353.in} (8 misspelt queries). Its variable-length,
hierarchical keys and duplicate records provide a check beyond the two synthetic
distributions.

\section{Results}

\subsection{Growth with dataset size}
\label{sec:growth}

For the linked list, node and string comparisons were exactly $N$ for every query in all
20 size--density combinations, confirming $\Theta(N)$ growth. Bit comparisons were also
linear in $N$ for fixed-length keys, but density changed the constant substantially:
Figure~\ref{fig:llbit} shows the dense-to-sparse ratio growing from about $5\times$ at
$N=50$ to about $10\times$ at $N=40\,000$, with wide sample-to-sample spread across
datasets and queries at the dense end (visible from the error bars). The linked list
stops comparing a stored key at
its first mismatching bit, so long pairwise agreement makes each of its $N$ comparisons
more expensive.

\begin{figure}[h]
\centering
\begin{tikzpicture}
\begin{semilogyaxis}[
  width=0.85\linewidth, height=5.0cm,
  xlabel={dataset size $N$}, ylabel={mean bit comparisons (log scale)},
  xtick={0,10000,20000,30000,40000}, scaled x ticks=false,
  xticklabels={0,10{,}000,20{,}000,30{,}000,40{,}000},
  legend style={at={(0.5,-0.32)}, anchor=north, legend columns=-1,
                font=\footnotesize, draw=none},
  grid=major, grid style={gray!20},
  error bars/y dir=both, error bars/y explicit,
]
\addplot[thick, blue, mark=*, mark size=1.2pt]
  table[x=n,y=ll_bit,y error=ll_bit_std,col sep=comma] {data/growth_sparse.csv};
\addlegendentry{Linked list, sparse}
\addplot[thick, blue, dashed, mark=square*, mark size=1.2pt]
  table[x=n,y=ll_bit,y error=ll_bit_std,col sep=comma] {data/growth_dense.csv};
\addlegendentry{Linked list, dense}
\addplot[thick, red, mark=*, mark size=1.2pt]
  table[x=n,y=pt_bit,y error=pt_bit_std,col sep=comma] {data/growth_sparse.csv};
\addlegendentry{Patricia trie, both densities}
\node[red!70!black, font=\scriptsize, anchor=south] at (axis cs:13000,264) {264 bits};
\end{semilogyaxis}
\end{tikzpicture}
\caption{Mean bit comparisons per exact query, with error bars showing $\pm 1$ standard
deviation across the pooled sample. Patricia results overlap at 264 bits with zero
variance; linked-list results grow linearly, with a density-dependent constant and
substantial dense-data spread.}
\label{fig:llbit}
\end{figure}

For the Patricia trie, every exact query compared exactly 264 bits: 32 characters plus
the 8-bit terminator, with zero variance across every dataset and query tested. This is
independent of both $N$ and density because the stems on a root-to-leaf path partition
the key bits, and it directly confirms the $\Theta(L)$ bound for a successful search (an
unsuccessful search only bounds this above, by $O(L)$, since it may stop before all $L$
bits are consumed).

Node accesses depend instead on how the distinguishing bits partition the keys. On sparse
data, the least-squares model $n_c=1.015\log_2 N+1.385$ follows the observations closely
(Figure~\ref{fig:ptnode}), consistent with approximately balanced branching and expected
$O(\log U)$ depth. Dense observations lie $2.6$--$4\times$ above this sparse-data fit over
the tested range, with substantial spread at every size. The fit is descriptive rather
than a formal bound: path length remains bounded by $O(\min\{U,L\})$, and a fixed $L$
ultimately caps tree depth regardless of density. The larger observed depth for dense
data is consistent with less even partitioning by the distinguishing bits, since Patricia
compression already stores the shared prefix itself in a single stem; partition balance
itself was not measured directly.

\begin{figure}[h]
\centering
\begin{tikzpicture}
\begin{axis}[
  width=0.85\linewidth, height=5.0cm,
  xlabel={dataset size $N$}, ylabel={mean node accesses},
  xtick={0,10000,20000,30000,40000}, scaled x ticks=false,
  xticklabels={0,10{,}000,20{,}000,30{,}000,40{,}000},
  legend style={at={(0.5,-0.32)}, anchor=north, legend columns=-1,
                font=\footnotesize, draw=none},
  grid=major, grid style={gray!20},
  error bars/y dir=both, error bars/y explicit,
]
\addplot[thick, red, mark=*, mark size=1.2pt]
  table[x=n,y=pt_node_sparse,y error=pt_node_sparse_std,col sep=comma] {data/growth_fit.csv};
\addlegendentry{Patricia, sparse}
\addplot[thick, black, dashed]
  table[x=n,y=log_fit,col sep=comma] {data/growth_fit.csv};
\addlegendentry{$1.015\log_2 N+1.385$}
\addplot[thick, red!60!black, mark=triangle*, mark size=1.6pt]
  table[x=n,y=pt_node_dense,y error=pt_node_dense_std,col sep=comma] {data/growth_fit.csv};
\addlegendentry{Patricia, dense}
\end{axis}
\end{tikzpicture}
\caption{Patricia node accesses, with error bars showing $\pm 1$ standard deviation
across the pooled sample. Sparse data track a logarithmic fit tightly; dense data have
deeper paths with much greater spread over the measured range.}
\label{fig:ptnode}
\end{figure}

At $N=40\,000$, relative to the linked list, the trie used $791\times$ fewer bit
comparisons and $2388\times$ fewer node accesses on sparse data. On dense data the
reductions were $7961\times$ and $599\times$, respectively. These ratios are specific to
this workload and should not be generalised as universal constants.

\FloatBarrier
\subsection{Edit position}
\label{sec:qtype}

Table~\ref{tab:qtype} fixes $N=5\,000$, pooling 5 datasets $\times$ 20 queries per row as
in Experiment 1. Linked-list node and string counts remain exactly 5000 with zero
variance, because the full scan is independent of query outcome. Patricia exact and
last-character-edit queries have similar bit and node counts because almost the whole key
is consumed before the decision. A first-character edit generally reduces bit
comparisons, but it does not guarantee an immediate trie mismatch: the modified prefix
may still follow a branch belonging to a different stored key, so the edit position and
the trie's actual mismatch position can differ.

This distinction is visible directly in the data. On sparse data, a first-character edit
reliably causes an early divergence (bit comparisons fall to $24.3\pm3.1$, a tight
distribution), because sparse keys branch almost immediately and few other keys share the
edited prefix. On dense data, the same edit is far less reliable: bit comparisons average
$101.8\pm56.7$, both much higher than the sparse case and far more variable, because 94\%
of dense keys still begin with the unedited baseline character, so a first-character edit
frequently lands on a prefix that other stored keys still share, and the search continues
until a real divergence is found deeper in the tree. The node effect follows the same
pattern: on sparse data a first-character edit changes the mean only from 13.7 to 12.7
nodes, well within one standard deviation, whereas on dense data it reduces the mean from
43.4 to 13.1 nodes, with a large standard deviation (18--20 nodes across all three dense
query types) that reflects exactly this edit-position/mismatch-position gap. Bit and node
comparisons therefore measure different things: bit comparisons track matched prefix
length, largely independent of tree shape, while node comparisons track the compressed
branching structure, which is far less stable for dense keys.

\begin{table}[h]
\centering\small
\begin{tabular}{llrrrr}
\toprule
Density & Query type & LL bit & LL node & PT bit & PT node \\
\midrule
Sparse & Exact                & $26195\pm655$      & 5000 & $264.0\pm0.0$  & $13.7\pm1.3$ \\
Sparse & Last-character edit  & $26168\pm640$      & 5000 & $252.9\pm1.1$  & $14.0\pm1.1$ \\
Sparse & First-character edit & $25992\pm614$      & 5000 & $24.3\pm3.1$   & $12.7\pm1.4$ \\
\midrule
Dense  & Exact                & $272285\pm119374$  & 5000 & $264.0\pm0.0$  & $43.4\pm17.9$ \\
Dense  & Last-character edit  & $236603\pm135990$  & 5000 & $253.8\pm0.8$  & $38.7\pm19.9$ \\
Dense  & First-character edit & $30109\pm3661$     & 5000 & $101.8\pm56.7$ & $13.1\pm5.5$ \\
\bottomrule
\end{tabular}
\caption{Mean $\pm$ standard deviation of comparisons per query at $N=5\,000$, pooled
across 5 datasets $\times$ 20 queries (100 samples) per row. LL = linked list; PT =
Patricia trie.}
\label{tab:qtype}
\end{table}

The low counter values for a first-character edit require careful interpretation even
when the edit does land near the tree's actual divergence point: a shallow mismatch can
still leave a large descendant subtree, so candidate enumeration and edit-distance
ranking may be expensive even though the counted search path is short. The table
therefore reflects only prefix-lookup cost; total spell-correction time is a separate,
unmeasured quantity.

\subsection{Provided dataset and duplicate keys}

The real dataset shows the same qualitative contrast (Table~\ref{tab:real}). The linked
list again performs exactly 1353 node and string comparisons per query. The Patricia trie
uses fewer than 20 node accesses on average, while its larger bit counts reflect the real
keys' greater and variable lengths.

\begin{table}[h]
\centering\small
\begin{tabular}{lrrrrrr}
\toprule
Query file & LL bit & LL node & LL str. & PT bit & PT node & PT str. \\
\midrule
\texttt{test1353.in} (588 exact)      & 116490 & 1353 & 1353 & 525.9 & 17.4 & 1.0 \\
\texttt{testpart1353.in} (8 misspelt) & 37536  & 1353 & 1353 & 328.9 & 9.1  & 1.0 \\
\bottomrule
\end{tabular}
\caption{Mean comparisons per query on \texttt{dataset\_1353.csv} ($N=1353$).}
\label{tab:real}
\end{table}

Duplicates affect the structures differently. In the linked list, each duplicate is
another record node and therefore increases the full-scan cost. In the Patricia trie,
records sharing a key are attached to one leaf, so multiplicity does not lengthen the
search path; it only increases the work needed to output the matching record list. This
separates $N$ from $U$: Patricia lookup cost depends primarily on the number of distinct
keys, independent of how many records each key returns.

\section{Discussion}

The measurements support the predictions in Table~\ref{tab:theory}. The linked list has
$\Theta(N)$ node and string comparisons for every query, with prefix density changing only
the constant in its bit count. Patricia bit comparisons are exactly bounded by key length
for a successful search ($\Theta(L)$, confirmed with zero variance) and bounded above for
an unsuccessful one ($O(L)$), while node depth is near logarithmic only when
distinguishing bits split the key set reasonably evenly. The dense experiment therefore
qualifies the average-case $O(\log U)$ claim while leaving Patricia's bit bound intact.

Whole-string comparisons illustrate why multiple counters are necessary: the linked list
performs one for every record, whereas this Patricia implementation performs at most one
after its bitwise descent. Bit and node counts expose different effects again in
Table~\ref{tab:qtype}. Reporting only one counter would conceal either prefix-density cost or
tree-shape cost.

With key length bounded, both structures use $\Theta(N)$ total space asymptotically: the
linked list allocates $\Theta(N)$ record nodes, and a binary Patricia trie with $U \le N$
leaves has at most $U-1$ internal nodes, so its total storage is $O(N+UL)$ in the worst
case and $O(N+U)=O(N)$ when $L$ is fixed. Patricia's practical disadvantage is therefore a
constant-factor one, not a different asymptotic class: each internal node additionally
stores a packed stem, a stem-length field, and two child pointers that the linked list
does not need. This implementation was not instrumented for bytes allocated, so this is
an asymptotic space observation rather than an empirical memory measurement.

Two limitations constrain interpretation. First, operation counts omit cache behaviour,
allocation overhead, output, candidate enumeration, and edit distance; they should not be
presented as wall-clock speedups. Second, each condition pools only five independently
generated datasets; this captures both dataset-to-dataset and query-to-query spread
(visible in Table~\ref{tab:qtype} and the error bars of Figures~\ref{fig:llbit}--\ref{fig:ptnode})
but is not a large sample, so the reported standard deviations themselves carry sampling
uncertainty, particularly for the highly variable dense condition. The systematic
variation in $N$, density, edit position, and real-data structure supports the growth
conclusions, but the exact ratios should not be generalised.

\section{Conclusion}

The experiments confirm the linked list's $\Theta(N)$ full-scan cost and the Patricia
trie's $\Theta(L)$ bitwise lookup for a successful search. Sparse keys produced
approximately logarithmic Patricia node depth, while dense keys produced deeper, more
variable paths, consistent with less even partitioning by their distinguishing bits. This
difference affects finite-sample node cost while leaving the bit bound intact. Overall,
the Patricia trie greatly reduces measured lookup work, especially as $N$ grows, while
duplicate handling, candidate ranking, and constant-factor memory overhead determine the
remaining trade-offs.

\end{document}
